% Generated draft: requires mathematical and source review.
\documentclass[11pt]{article}
\usepackage[margin=25mm]{geometry}
\usepackage{fontspec}
\setmainfont{Latin Modern Roman}
\usepackage{amsmath,amssymb,mathtools,unicode-math}
\setmathfont{Latin Modern Math}
\usepackage{enumitem,needspace,xurl,hyperref}
\setlength{\parindent}{0pt}
\setlength{\parskip}{5pt}
\setlength{\emergencystretch}{3em}
\newcommand{\N}{\mathbb{N}}
\newcommand{\Z}{\mathbb{Z}}
\newcommand{\Q}{\mathbb{Q}}
\newcommand{\R}{\mathbb{R}}
\newcommand{\C}{\mathbb{C}}
\newcommand{\F}{\mathbb{F}}
\newcommand{\A}{\mathbb{A}}
\newcommand{\PP}{\mathbb{P}}
\newcommand{\E}{\mathbb{E}}
\newcommand{\eps}{\varepsilon}
\newcommand{\dd}{\,\mathrm{d}}
\begin{document}
\section*{UnsolvedMath ID: 10600060}
\textbf{Source problem code:} catalogue entry
\subsection{Definitions and mathematical statement}
A free knot is represented by a framed \(4\)-valent graph with one unicursal component, modulo the intended free-knot Reidemeister moves; framing specifies opposite half-edges. For the source's spanning-cobordism category, define the slice genus \(g_s(K)\) as the minimum genus of an admissible spanning surface with boundary \(K\), or \(+\infty\) if none exists. The problem is to extend sliceness-obstruction methods to lower bounds on \(g_s(K)\) for arbitrary-genus cobordisms, and investigate the stronger parity/group invariants described below.

The source gives two obstacles. First, the path joining two preimages of a generic double-line point need not be homotopically unique on a positive-genus spanning surface, so parity and justified parity need not be well defined. It proposes vanishing of the \(\Z/2\Z\)-class dual to the double-line graph \(\Psi\), noting a connection with atoms and even-valent graphs. Second, the Morse/Reeb structure can allow splitting and recombination that cancels signed invariant values: its schematic example starts with \(L(K)=8\), splits into two components with \(L(K_i)=4\), and then allows cancellation \(4+(-4)=0\). This is presented as an obstruction to the proposed method, not a proved construction of a specific movie. The phrase that a Reeb graph is not necessarily a disc requires clarification, since a graph and a two-dimensional disc are different objects.

For a spanning surface \(\mathcal D_g\) with boundary \(S^1\), the source assigns the double-line collection a class
\[
 \kappa\in H_1(\mathcal D_g,S^1;\Z/2\Z).
\]
It reports that this is the sole obstruction to the path-independent even/odd double-line parity and to checkerboard colouring, and defines the cobordism to be \emph{checkerboard} or \emph{atomic} when \(\kappa=0\). Within this class, to distinguish the source's odd labels \(b,b'\), one joins the two preimages of an odd-stratum point and counts intersections with even double lines. The remaining obstruction is the class of the even double lines,
\[
 \kappa'\in H_1(\mathcal D_g,S^1;\Z/2\Z).
\]
An atomic cobordism is \emph{\(2\)-atomic} when also \(\kappa'=0\). The source states the theorem that if a one-component framed \(4\)-graph \(K\) satisfies \(L(K)\ne0\), then no \(2\)-atomic cobordism of \emph{any} genus spans \(K\). This is a reported theorem only for the restricted class, not an unrestricted slice-genus bound.

The continuation describes a strengthening \(G_m\) of the group \(G=G_1\), with invariants valued in conjugacy classes of \(G_m\), obtained by refining even chords into types and related to iterated parity and deletion of odd crossings. It proposes that all these invariants obstruct sliceness, and asks to investigate their relation to cobordisms. It also reports virtual-knot variants retaining over/under information in addition to parity. The definitions of \(L,G,G_m\), the generators/relations, admissible surface singularities, parity counts and cited proofs are missing; neither their full sliceness property nor any unrestricted genus estimate is asserted here.
\subsection{Short English statement}
Seek slice-genus bounds for free knots beyond genus-zero arguments, accounting for double-line homology obstructions and cancellation under Morse moves, and determine the cobordism strength of the refined group-valued parity invariants.
\subsection{Sources}
\begin{itemize}
\item \emph{Unsolved Problems in Virtual Knot Theory and Combinatorial Knot Theory}, revised survey, problem-list item 60, ``Problems on Free Knot Cobordism.'' \url{https://arxiv.org/abs/1409.2823}. Recorded access 2026-07-29; extraction source-tex.
\item The inline references for \(L\), atom connections, the \(2\)-atomic theorem, the \(G_m\) papers, virtual-knot variants and a separate cobordism paper are absent. The text attributes the group strengthening to ``the fourth author,'' whose identity is not supplied; it also refers to missing section and figure labels. No author or bibliography entry is reconstructed.
\item Local UnsolvedMath registry, \texttt{data/unsolvedmath/problems.json}, permanent ID \(10600060\), exact source code \texttt{catalogue entry}; supplied partially\_solved/PARTIAL-PROGRESS status. Its summary distinguishes established sliceness-obstruction work from the continuing arbitrary-genus lower-bound problem.
\end{itemize}
\end{document}
